% RESEARCH_DRAFT_SHA256: 339b585f2f0d870448e1b7f857d550eea647f76503c62df17c2cf5767e1fd772
% Layout and source-reference wrappers expanded for the web renderer.
\documentclass[11pt]{article}
\usepackage[margin=25mm]{geometry}
\usepackage{fontspec,unicode-math}
\setmainfont{Latin Modern Roman}
\setmathfont{Latin Modern Math}
\usepackage{amsmath,amssymb,mathtools,xurl,hyperref}
\hypersetup{colorlinks=true,urlcolor=blue,linkcolor=blue}
\setlength{\parindent}{0pt}
\setlength{\parskip}{5pt}
\setlength{\emergencystretch}{3em}
\newcommand{\sref}[2]{\hyperlink{source:#1:#2}{[S#2]}}
\newcommand{\cataloguescope}[1]{\paragraph{Source question.}#1}
\newcommand{\floor}[1]{\left\lfloor#1\right\rfloor}
\begin{document}

\section*{Repeated integer-infeasibility question: a sharp exponential family that fails fixed determinants}
\noindent catalogue entry; catalogue entry. Research attempt dated 11 October 2026.
\paragraph{Result in scope.}This manuscript contains a source-specific attempt and complete proofs of the stated auxiliary calculations. It makes no new solution or novelty claim for the entire target.
\section{Original problem and preserved definition}
The following definition and source list are reproduced from the repository. Any earlier inspection dates and reported literature claims in this preserved text describe that source record, not additional verification by this attempt.
% BEGIN PRESERVED DEFINITION
\subsection{Definitions and mathematical statement}

\paragraph{Shared notation.}Standard number systems are $\mathbb N,\mathbb Z,\mathbb Q,\mathbb R,\mathbb C$. Specialized notation and index ranges follow the cited source and the definitions below.


The constraint matrix $A$ is integral, with rows defining the displayed inequalities. Integer-infeasible means $P\cap\mathbb Z^n=\varnothing$; minimally integer-infeasible means deleting any one inequality produces a set containing an integer point. Full-dimensional means that $P$ has nonempty interior in $\mathbb R^n$. The parameter $\Delta$ is the maximum absolute determinant among all square submatrices of $A$.
\paragraph{Record scope and qualifications.}
The catalogue title is retained as metadata. Named certificates, contracts, bounds or frameworks in that title are not asserted as results of the original source extract.
The supplied registry flags this extraction as unrecoverable. It is displayed as a literal source fragment; the definitions provide context rather than a replacement formal statement.

Source-record limitation: Source flags the statement as unrecoverable: The official report confirms that this is a duplicate extraction of the minimally integer-infeasible systems problem retained in OWR-8415356-014. No separate clean mathematical statement is supplied; this source ID preserves its provenance.

Related catalogue records: \href{https://mehmetmars7.github.io/Erdosproblems-llm-hunter/problem.html?type=open_problems&id=30004903}{30004903}. Separate source IDs are preserved for their provenance.
The source text below is rendered literally where necessary; damaged formulas and composite extracts are not silently repaired into a new mathematical claim.
The supplied catalogue marks this record as unrecoverable. Its extract and terminology are documented, but a complete original problem has not been reconstructed.
\paragraph{Statement recorded in the supplied source.}
Can we upper bound m as a linear function of n, provided that ∆ is con- stant? Open Problem: Integer infeasible subsystems Robert Weismantel (joint work with Marcel Celaya) \href{https://doi.org/10.4171/owr/2021/53}{[S2]}
\paragraph{Source status.}
The supplied catalogue labels this source record open; that metadata was checked against the supplied snapshot on 2026-10-11. The label is preserved for provenance and is not an independent verification of present open status or mathematical correctness. The catalogue literature assessment is dated 2026-08-22.
\subsection{Short English statement}
The supplied fragment concerns linear Bounds for Minimally Integer-Infeasible Systems. Its complete mathematical target is not recovered.
\subsection{Sources}
\begin{itemize}
\item[S1] Ulam AI and the UnsolvedMath Contributors, UnsolvedMath: A Curated Collection of Open Mathematics Problems, catalogue entry. Catalogue curation and metadata: CC BY 4.0. The source extract is retained; definitions, English summaries and qualifications are editorial additions.
\url{https://huggingface.co/datasets/ulamai/UnsolvedMath}.
\item[S2] Combinatorial Optimization (2022), Oberwolfach Reports, DOI 10.4171/owr/2021/53, pp. 2893--2954
\url{https://doi.org/10.4171/owr/2021/53}.
\item[S3] JCTA bounded-subdeterminant bounds
\url{https://doi.org/10.1016/j.jcta.2024.106003}.
\item Source-specific Oberwolfach Reports licences, including CC BY-SA 4.0 where applicable, are retained; see THIRD\_PARTY\_NOTICES.md.
\end{itemize}
% END PRESERVED DEFINITION
\subsection{Verbatim supplied catalogue question}
Can we upper bound m as a linear function of n, provided that ∆ is con- stant? Open Problem: Integer infeasible subsystems Robert Weismantel (joint work with Marcel Celaya)

\section{Duplicate scope and attempted counterexample family}
This is the trailing duplicate of the minimally integer-infeasible
systems question in record 30004903. The original context on printed
page 2943 supplies the omitted matrix, integrality and full-dimensionality
hypotheses. We test a natural exponential-facet construction against
the crucial fixed-subdeterminant condition. The family has all the
other required properties but fails that condition, so it does not
disprove the source question.

\paragraph{An exact minimally integer-infeasible system.}
For $n\ge2$ and every sign vector $\sigma\in\{-1,1\}^n$,
include the inequality
\[
2\sigma\cdot x\le2p(\sigma)-1,
\qquad p(\sigma)=|\{i:\sigma_i=1\}|.
\]
There are $2^n$ integral rows. Writing
$c=(1/2,\ldots,1/2)$, the inequalities become
$\sigma\cdot(x-c)\le(n-1)/2$ for every sign vector.
The maximum of $\sigma\cdot y$ over all signs is
$\sum_i|y_i|$, obtained by choosing the sign of each
coordinate independently. Thus their intersection is exactly
\[
P=\{x:\sum_i|x_i-1/2|\le(n-1)/2\}.
\]
It has nonempty interior because the center satisfies every
inequality strictly and the radius is positive. Every integer
coordinate has distance at least $1/2$ from $1/2$, so
every integer vector has total distance at least $n/2$ and
lies outside $P$.

For each row $\sigma$, let $z_\sigma$ have coordinate one
where $\sigma_i=1$ and zero otherwise. It violates its own
row, since $2\sigma\cdot z_\sigma=2p(\sigma)$.
For another sign vector $\tau$, let $h$ be the number of
positions where it differs from $\sigma$. Directly counting
positive and negative positions gives
$\tau\cdot z_\sigma=p(\tau)-h$. If $\tau\ne\sigma$
then $h\ge1$, whence
$2\tau\cdot z_\sigma\le2p(\tau)-2<2p(\tau)-1$.
Thus deleting row $\sigma$ admits the integer point
$z_\sigma$. Every row is essential and the system is minimally
integer-infeasible, with $m=2^n$ and entries all in $\{-2,2\}$.

\paragraph{The determinant promise destroys this counterexample route.}
Bounded entries are not bounded subdeterminants. For $n=2^s$,
define the sign matrix $H_1=(1)$ and recursively
\[
H_{2r}=\begin{pmatrix}H_r&H_r\\H_r&-H_r\end{pmatrix}.
\]
Block multiplication shows inductively that
$H_nH_n^T=nI$. Taking determinants gives
$|\det H_n|=n^{n/2}$. Its rows are sign vectors and are
distinct, since nonzero orthogonal rows cannot coincide. Consequently
the full constraint matrix contains the square row submatrix $2H_n$,
whose determinant has magnitude $2^n n^{n/2}$.
Therefore its maximum subdeterminant $\Delta$ tends to infinity
along these dimensions, even though every entry has magnitude two.

The exponential number of essential inequalities proves only that
one cannot replace the source's bound on every minor by a bound
on individual coefficients. The inequality $m\le nf(\Delta)+g(\Delta)$
allows arbitrary growth in $\Delta$, so this family with growing
$\Delta$ contradicts no such bound. A valid negative construction
would need infinitely many dimensions with one common determinant
bound; no such construction is obtained here.

\section{Audit}
The original source paragraph and duplicate provenance were inspected.
The audit checked full-dimensionality, integrality of right sides,
every deleted-row witness and the actual large minor, not just
entry bounds. All arguments are exact hand calculations; no determinant
enumeration was run. This is a classical-style sharp parity family
used to diagnose a failed counterexample route, not a new independent
problem or a disproof of the original fixed-$\Delta$ question.

\section{Process and contributions}
Mathematics and model-assisted adversarial self-audit: OpenAI gpt-6-astra, ultra reasoning, as dispatched for this research batch; exact serving revision not exposed. Codex (GPT-6; exact serving revision not exposed) materially assisted with orchestration and format review. No human mathematical authorship, mathematical edits, or independent human proof checking is asserted.
Contribution prompt A was applied to the full definition and supplied source row, then prompt B was applied adversarially to the derived results. The item-specific audit above records the scope checks and rejected extrapolations. Local assembly and structural checks do not establish mathematical correctness. The source catalogue status is preserved in TOP\_PROBLEM and is separate from this attempt status.
\paragraph{Integration review.} Codex (GPT-6; exact serving revision not exposed) reviewed the written arguments and expanded layout and source-reference wrappers. Repository metadata and the production web renderer were checked using Python 3.13 and JavaScript with MathJax 3.2.2. This model-assisted review is not independent human proof verification; no PDF compilation or formal proof checking is asserted.

\section{Completion Estimate}
\noindent\textbf{COMPLETION: 0\%}
The proved material is a rederivation of known or auxiliary facts rather than a substantiated new advance on the original target. The exponential essential-inequality family has unbounded higher minors, so it cannot resolve the original fixed-Delta linear-bound question.
\section{References}
The complete original Sources subsection above is retained, including catalogue attribution and licence notices. The additional list records primary sources actually consulted for this attempt on 11 October 2026; uncited source links reproduced above were not all rechecked.
\begin{itemize}
\item Oberwolfach Report 53/2021, Combinatorial Optimization, original question and surrounding discussion read from the primary PDF text on 11 October 2026. Historical statements are not a current-status survey. \url{https://ems.press/content/serial-article-files/46931}.
\end{itemize}
\section{Final status}
\textbf{Unresolved in this attempt.} The exponential essential-inequality family has unbounded higher minors, so it cannot resolve the original fixed-Delta linear-bound question.
\end{document}
